\section{Related Work}
\label{sec:related}

Bodeveix and Filali study type synthesis and the translation of B to
PVS~\cite{Bodeveix2002BtoPVS}. Their work addresses the passage from B's
set-based language to the typing discipline of the target prover.
\barrel shares this concern and targets obligations exported by Atelier~B,
with an interactive interface hosted in Lean. Its mathematical objects
use Mathlib sets and relations, allowing the imported propositions to be
proved with the host prover's tactics and libraries.

Proof requirements attached to partially specified operations have
established precedents. Sozeau's Russell language adapts predicate
subtyping from PVS to Coq, generating obligations whose proofs complete
a well-typed term~\cite{Sozeau2007SubsetCoercions}. Schmalz studies
partial functions, definitional extensions, and reasoning about the
logic of Event-B, with an Isabelle/HOL integration for
Rodin~\cite{Schmalz2012RodinIsabelle}. These approaches provide relevant
foundations for expressing partiality in a total logic. \barrel applies
explicit proof requirements at occurrences of supported B operators,
closes the resulting conditions over their local context, and connects
them to automatic and interactive discharge of Atelier~B obligations.

The Z Mathematical Toolkit in Isabelle/HOL develops a Z metalanguage
and type universe over HOL, connected to concrete representations such
as finite functions and lists. Coercive subtyping and overloading support
Z-style sets, relations, functions, and sequences~\cite{Foster2025ZToolkit}.
\barrel's library serves the corresponding role for the B fragment
accepted by its translator in Lean. The availability of reusable
mathematical lemmas is important in both settings: the usefulness of
an embedding depends on how users can reason about its representations,
as well as on the expressions it can translate.

Ballenghien and Wolff's HOL-B prototype hosts Event-B as a domain-specific
language in Isabelle/HOL, with a mathematical toolkit, IDE and
documentation support, and generated denotational semantics for machine
specifications~\cite{Wolff2024EB_Isabelle}. This places machine semantics
inside the proof assistant. \barrel retains Atelier~B's proof-obligation
generation and consumes the exported formulas. The resulting workflow
fits existing B developments, while leaving the correctness of the
external generator and of the translation outside the kernel guarantee.
This architectural choice is distinct from interpreting entire machines
inside the prover.

Automated backends complement the interactive approach. D\'eharbe studies
the integration of SMT solvers with B and Event-B environments~\cite{DeharbeBSMT},
and BEer translates B proof obligations into SMT-LIB~\cite{beer}.
\barrel instead exposes propositions and proof tasks in Lean: users may
combine automatic tactics with structured interactive arguments, prove
auxiliary lemmas, and inspect the resulting dependencies. Its automation
is intended to handle routine conditions within that proof workflow.
The case study examines this division of work through the
WD obligations discharged automatically and the principal proofs that
remain, rather than a competition between solver success rates.
